\section{Risultati sperimentali}
\label{sec:risultati}

\subsection{Confronto fra le varianti del modello}

La seguente tabella riporta i risultati del confronto
sperimentale fra il modello originale e le tre varianti compresse, misurati
sul test set corretto dalla suite di benchmark integrata nella pipeline di
compressione (\texttt{models/compress.py}), che salva i valori in
\texttt{benchmarks/benchmark\_results.csv}.

\begin{table}[h]
\centering
\small
\begin{tabular}{@{}lccc@{}}
\toprule
\textbf{Variante} & \textbf{Accuracy} & \textbf{Dimensione} & \textbf{Latenza media} \\
\midrule
Baseline (float32) & 95,0\% & 33,44 MB & 3,33 ms \\
Pruned (sparsità 30\%) & 82,2\% & 33,44 MB & 3,28 ms \\
Pruned + Quantized (INT8) & 82,2\% & 9,43 MB & 4,32 ms \\
ONNX Quantized (INT8) & 95,0\% & 8,37 MB & 6,10 ms \\
\bottomrule
\end{tabular}
\caption{Confronto sperimentale fra baseline e varianti compresse, sul test set corretto.}
\label{tab:benchmark-finale}
\end{table}

\subsection{Interpretazione dei risultati}

I risultati permettono di trarre alcune conclusioni rilevanti sul
trade-off fra le diverse tecniche di compressione applicate a questa
specifica architettura:

\begin{itemize}
    \item \textbf{La sola quantizzazione (variante ONNX) preserva quasi
    completamente l'accuratezza}, riducendo comunque la dimensione del
    modello di un fattore $\sim$4. Questo risultato è coerente con la
    composizione dell'architettura: \texttt{TomatoCNN} è dominata da strati
    convoluzionali (\texttt{Conv2d}), sui quali la quantizzazione dinamica
    ha un impatto contenuto sulla precisione numerica rispetto a quanto
    avverrebbe su architetture dominate da strati completamente connessi.
    \item \textbf{Il pruning strutturato al 30\% introduce una perdita di
    accuratezza significativa} (da 95,0\% a 82,2\%), a fronte di
    \emph{nessuna} riduzione di dimensione su disco (33,44 MB, identica alla
    baseline). Il motivo è che il pruning si limita ad azzerare i pesi dei
    filtri scartati, ma il file salvato su disco continua a riservare lo
    stesso spazio per ciascun peso, zero compreso: un filtro azzerato pesa
    esattamente come uno non azzerato, finché non lo si salva in un formato
    che memorizza esplicitamente solo i valori diversi da zero (un formato
    "sparso"), operazione che la pipeline attuale non esegue. Il pruning da
    solo, quindi, non produce alcun beneficio pratico in termini di
    spazio occupato: serve sempre in combinazione con la quantizzazione (o
    con un formato di salvataggio sparso) per tradursi in un file più
    piccolo.
    \item \textbf{La combinazione pruning + quantizzazione} produce il file
    più compatto fra le due varianti PyTorch (9,43 MB), ma eredita
    interamente la perdita di accuratezza introdotta dal pruning
    (82,2\%), risultando la variante col trade-off meno favorevole fra le
    tre, a meno che il vincolo di memoria del dispositivo target non sia
    particolarmente stringente.
    \item Le \textbf{latenze misurate} risultano tutte dello stesso ordine
    di grandezza (3--6 millisecondi per singola inferenza), un dato atteso
    data la ridotta dimensione della rete; le differenze osservate fra le
    varianti non sono quindi determinanti nella scelta della configurazione
    da adottare per questo specifico modello, mentre lo sarebbero per
    architetture più profonde e con un numero di parametri maggiore.
\end{itemize}

In sintesi, per questa specifica architettura e questo specifico compito di
classificazione, \textbf{la quantizzazione via ONNX Runtime rappresenta la
scelta più equilibrata}, offrendo una riduzione di dimensione significativa
a fronte di una perdita di accuratezza trascurabile; il pruning strutturato,
pur essendo una tecnica rilevante in architetture più grandi o con maggiore
ridondanza nei filtri, in questo caso specifico penalizza l'accuratezza in
modo più marcato di quanto benefici in termini di prestazioni.

\subsection{Metriche di comportamento dell'agente}

Oltre alle metriche puramente relative al modello di classificazione, il
sistema misura anche il comportamento dell'agente decisionale nel suo
complesso: numero di cicli eseguiti, numero di rilevamenti di malattia
consecutivi, e \emph{behavior accuracy}\footnote{La \emph{behavior
accuracy} è la percentuale di cicli in cui l'agente si è comportato in
modo appropriato rispetto alla reale condizione della pianta: quanto
spesso, cioè, la sua reazione (attivare la ventilazione, l'irrigazione,
un allarme, o non fare nulla) è risultata quella corretta per la
situazione realmente presente, non simulata. Non coincide quindi con
l'accuratezza pura della CNN: un ciclo è considerato corretto non solo
quando la classe predetta coincide con quella reale, ma anche quando
entrambe appartengono alla categoria delle malattie fungine (Early
Blight, Late Blight, Leaf Mold, Septoria Leaf Spot, Target Spot), pur
essendo classi diverse. La ragione è che, per tutte le malattie fungine,
l'agente applica la stessa regola decisionale (attivazione della
ventilazione in presenza di umidità alta): se la CNN scambia una fungina
per un'altra, la predizione "pura" è sbagliata, ma l'azione
effettivamente intrapresa dall'agente è comunque quella corretta per la
situazione reale. La metrica misura quindi quanto spesso l'agente si è
comportato correttamente, non quanto spesso la CNN ha indovinato
l'etichetta esatta -- ed è per questo tipicamente più alta
dell'accuratezza di classificazione riportata in
Tabella~\ref{tab:benchmark-finale}.} (la frazione di cicli in cui l'azione
dell'agente è risultata coerente con la classe reale nota dal dataset). Le
prime due
sono visibili sia nel log CSV prodotto dall'orchestratore da riga di
comando sia nella dashboard interattiva; la \emph{behavior accuracy} è
invece calcolata e mostrata solo nel riepilogo finale dell'orchestratore da
riga di comando, al termine della sessione. Queste metriche permettono di
valutare non solo la qualità della singola predizione della CNN, ma la
correttezza del comportamento del sistema nel suo complesso (ad esempio, se
l'attivazione della ventilazione sia realmente avvenuta in presenza di
condizioni favorevoli a una patologia fungina).
